\chapter{Introducere}
\label{ch:intro}

\section{Context și motivație}

Aplicațiile web moderne servesc un volum mare de cereri de citire asupra unor date relativ
stabile: pagini de produs, profile, clasamente, configurații. Baza de date relațională rămâne sursa
de adevăr datorită garanțiilor tranzacționale (ACID), dar fiecare interogare implică rețea,
analiză sintactică, planificare, acces la pagini și eventual agregări costisitoare. Atunci când
aceeași informație este cerută de mii de ori pe secundă, recalcularea ei la fiecare cerere este
risipă de resurse și, mai ales, sursă de latență.

Un \emph{cache in-memory} păstrează rezultate deja calculate în memoria RAM, aproape de aplicație,
astfel încât cererile repetate să fie servite în microsecunde, fără a atinge baza de date. Practic
toate sistemele web la scară mare folosesc o astfel de soluție (Memcached la Facebook
\cite{memcache-fb}, Redis, cache-uri locale de tip Caffeine/Guava etc.). Problema este însă mai
subtilă decât pare: memoria este limitată (trebuie ales \emph{ce} păstrăm), datele se modifică
(trebuie menținută coerența), iar sistemul este distribuit (mai multe instanțe, mai multe noduri).

\section{Obiectivele proiectului}

Proiectul urmărește proiectarea și realizarea unei aplicații paralele și distribuite care să răspundă
concret acestor probleme:

\begin{enumerate}
  \item un \textbf{motor de tabele in-memory} (\texttt{memtable}) thread-safe, cu sharding, politici de
        evacuare LRU/FIFO/CLOCK, TTL și invalidare pe grupuri de chei (tag-uri);
  \item un \textbf{cache distribuit}: noduri de cache accesate printr-un protocol TCP propriu și un client
        care distribuie cheile prin \emph{consistent hashing};
  \item o \textbf{aplicație web} realistă (catalog de produse cu recenzii în PostgreSQL) care folosește
        cache-ul prin mai multe strategii de citire și scriere, comutabile la runtime;
  \item \textbf{comparația cu alternative}: Redis și o tabelă in-memory relațională (tabelă
        \texttt{UNLOGGED} într-un PostgreSQL cu date pe \texttt{tmpfs});
  \item o \textbf{evaluare experimentală} riguroasă: debit, latențe pe percentile, rata de hit,
        scalabilitate (fire, clienți, instanțe, noduri) și comparația cu varianta secvențială.
\end{enumerate}

\section{Contribuții}

\begin{itemize}
  \item Implementare completă în Rust (≈3.400 de linii), fără cod \texttt{unsafe}, verificată cu
        \texttt{clippy} în modul \emph{pedantic} și acoperită de teste unitare și de integrare.
  \item Infrastructură reproductibilă cu Docker Compose: 12 containere, observabilitate cu
        Prometheus și Grafana, documentație OpenAPI (Swagger) și o interfață web de control.
  \item Mecanisme avansate: \emph{single-flight} (protecție la \emph{cache stampede}), cache pe două
        niveluri (\emph{near-cache}), invalidare distribuită prin \texttt{LISTEN/NOTIFY}, write-behind cu
        grupare în loturi și oprire fără pierderi de date.
  \item Un generator de încărcare propriu cu distribuție Zipf și histograme HDR și un simulator
        trace-driven pentru politicile de evacuare, comparate cu limita teoretică.
\end{itemize}

\section{Structura documentului}

\Cref{ch:fundamente} prezintă fundamentele teoretice. \Cref{ch:analiza} definește problema și
cerințele și compară arhitectura secvențială cu cea paralelă/distribuită. \Cref{ch:arhitectura}
descrie arhitectura sistemului, iar \cref{ch:implementare} detaliile de implementare.
\Cref{ch:tehnologii} trece în revistă tehnologiile utilizate. \Cref{ch:evaluare} conține
metodologia și rezultatele experimentale, iar \cref{ch:concluzii} concluziile. Anexa
\ref{ch:rulare} descrie rularea proiectului.
